\documentclass[11pt,letterpaper]{article}

\usepackage[margin=1in]{geometry}
\usepackage{amsmath}                  % math in the GNN section
\usepackage{booktabs}                % rules in the data tables
\usepackage{longtable}               % multi-page per-target tables
\usepackage[labelfont=bf]{caption}   % bold "Table SN" labels

% renumber figures/tables/equations as S1, S2, ...
\renewcommand{\thetable}{S\arabic{table}}
\renewcommand{\thefigure}{S\arabic{figure}}
\renewcommand{\theequation}{S\arabic{equation}}

% per-target table bodies, auto-generated by src/make_si_tables.py
% auto-generated by src/make_si_tables.py -- do not edit by hand
\newcommand{\siCorrRows}{%
CHEMBL1862\_Ki & 161 & 41.6 & 0.62 & 0.38 & 0.33 & 0.23 & 0.25 \\
CHEMBL1871\_Ki & 134 & 23.9 & 0.76 & 0.51 & 0.32 & 0.21 & 0.25 \\
CHEMBL2034\_Ki & 152 & 34.2 & 0.77 & 0.65 & 0.45 & 0.34 & 0.35 \\
CHEMBL2047\_EC50 & 128 & 39.1 & 0.69 & 0.51 & 0.34 & 0.28 & 0.29 \\
CHEMBL204\_Ki & 553 & 39.8 & 0.64 & 0.53 & 0.40 & 0.30 & 0.41 \\
CHEMBL2147\_Ki & 294 & 39.5 & 0.50 & 0.40 & 0.19 & 0.15 & 0.16 \\
CHEMBL214\_Ki & 666 & 36.8 & 0.65 & 0.43 & 0.27 & 0.16 & 0.22 \\
CHEMBL218\_EC50 & 208 & 36.5 & 0.69 & 0.44 & 0.30 & 0.16 & 0.26 \\
CHEMBL219\_Ki & 374 & 39.6 & 0.70 & 0.50 & 0.35 & 0.15 & 0.29 \\
CHEMBL228\_Ki & 342 & 37.1 & 0.65 & 0.45 & 0.25 & 0.12 & 0.19 \\
CHEMBL231\_Ki & 197 & 24.4 & 0.67 & 0.45 & 0.32 & 0.22 & 0.21 \\
CHEMBL233\_Ki & 630 & 41.1 & 0.69 & 0.47 & 0.33 & 0.22 & 0.29 \\
CHEMBL234\_Ki & 733 & 43.7 & 0.63 & 0.43 & 0.28 & 0.13 & 0.18 \\
CHEMBL235\_EC50 & 472 & 37.7 & 0.66 & 0.43 & 0.26 & 0.16 & 0.24 \\
CHEMBL236\_Ki & 521 & 39.2 & 0.64 & 0.40 & 0.28 & 0.21 & 0.27 \\
CHEMBL237\_EC50 & 193 & 47.2 & 0.63 & 0.31 & 0.22 & 0.16 & 0.20 \\
CHEMBL237\_Ki & 522 & 42.7 & 0.67 & 0.39 & 0.27 & 0.14 & 0.26 \\
CHEMBL238\_Ki & 213 & 25.8 & 0.63 & 0.45 & 0.21 & 0.21 & 0.19 \\
CHEMBL239\_EC50 & 347 & 40.6 & 0.65 & 0.42 & 0.30 & 0.16 & 0.27 \\
CHEMBL244\_Ki & 621 & 47.7 & 0.63 & 0.51 & 0.34 & 0.25 & 0.31 \\
CHEMBL262\_Ki & 173 & 18.5 & 0.63 & 0.39 & 0.12 & 0.11 & 0.13 \\
CHEMBL264\_Ki & 574 & 41.6 & 0.62 & 0.37 & 0.27 & 0.17 & 0.20 \\
CHEMBL2835\_Ki & 126 & 10.3 & 0.74 & 0.59 & 0.49 & 0.36 & 0.20 \\
CHEMBL287\_Ki & 267 & 38.6 & 0.61 & 0.38 & 0.19 & 0.01 & 0.15 \\
CHEMBL2971\_Ki & 197 & 17.3 & 0.73 & 0.52 & 0.54 & 0.32 & 0.24 \\
CHEMBL3979\_EC50 & 226 & 41.6 & 0.60 & 0.41 & 0.10 & 0.02 & 0.05 \\
CHEMBL4005\_Ki & 193 & 42.0 & 0.71 & 0.63 & 0.13 & 0.14 & 0.06 \\
CHEMBL4203\_Ki & 149 & 8.7 & 0.62 & 0.61 & 0.19 & 0.11 & 0.16 \\
CHEMBL4616\_EC50 & 139 & 52.5 & 0.73 & 0.48 & 0.15 & 0.02 & 0.13 \\
CHEMBL4792\_Ki & 297 & 53.9 & 0.72 & 0.47 & 0.23 & 0.16 & 0.24 \\
}

\newcommand{\siGnnRows}{%
CHEMBL1862\_Ki & 161 & 0.59 & 0.78 & 0.88 & 0.049 \\
CHEMBL1871\_Ki & 134 & 0.53 & 0.61 & 0.71 & 0.056 \\
CHEMBL2034\_Ki & 152 & 0.50 & 0.59 & 0.83 & 0.074 \\
CHEMBL2047\_EC50 & 128 & 0.42 & 0.64 & 0.69 & 0.078 \\
CHEMBL204\_Ki & 553 & 0.55 & 0.80 & 0.74 & 0.058 \\
CHEMBL2147\_Ki & 294 & 0.47 & 1.14 & 0.92 & 0.053 \\
CHEMBL214\_Ki & 666 & 0.51 & 0.71 & 0.69 & 0.063 \\
CHEMBL218\_EC50 & 208 & 0.53 & 0.69 & 0.80 & 0.064 \\
CHEMBL219\_Ki & 374 & 0.55 & 0.77 & 0.77 & 0.073 \\
CHEMBL228\_Ki & 342 & 0.52 & 0.67 & 0.74 & 0.062 \\
CHEMBL231\_Ki & 197 & 0.56 & 0.70 & 0.68 & 0.051 \\
CHEMBL233\_Ki & 630 & 0.60 & 0.76 & 0.73 & 0.065 \\
CHEMBL234\_Ki & 733 & 0.49 & 0.66 & 0.72 & 0.065 \\
CHEMBL235\_EC50 & 472 & 0.48 & 0.61 & 0.68 & 0.072 \\
CHEMBL236\_Ki & 521 & 0.52 & 0.81 & 0.69 & 0.066 \\
CHEMBL237\_EC50 & 193 & 0.58 & 0.83 & 0.86 & 0.089 \\
CHEMBL237\_Ki & 522 & 0.55 & 0.70 & 0.67 & 0.068 \\
CHEMBL238\_Ki & 213 & 0.49 & 0.63 & 0.68 & 0.061 \\
CHEMBL239\_EC50 & 347 & 0.51 & 0.68 & 0.94 & 0.085 \\
CHEMBL244\_Ki & 621 & 0.56 & 0.80 & 0.76 & 0.069 \\
CHEMBL262\_Ki & 173 & 0.57 & 0.75 & 0.76 & 0.063 \\
CHEMBL264\_Ki & 574 & 0.48 & 0.62 & 0.66 & 0.066 \\
CHEMBL2835\_Ki & 126 & 0.27 & 0.33 & 0.41 & 0.063 \\
CHEMBL287\_Ki & 267 & 0.56 & 0.66 & 0.75 & 0.071 \\
CHEMBL2971\_Ki & 197 & 0.43 & 0.58 & 0.67 & 0.043 \\
CHEMBL3979\_EC50 & 226 & 0.49 & 0.66 & 0.69 & 0.076 \\
CHEMBL4005\_Ki & 193 & 0.50 & 0.66 & 0.87 & 0.072 \\
CHEMBL4203\_Ki & 149 & 0.67 & 0.79 & 0.77 & 0.052 \\
CHEMBL4616\_EC50 & 139 & 0.51 & 0.65 & 0.76 & 0.110 \\
CHEMBL4792\_Ki & 297 & 0.54 & 0.68 & 0.73 & 0.096 \\
}


\title{Supporting Information\\[4pt]
\large Structure--Activity Landscape Roughness Predicts Per-Compound QSAR Error Independently of the Applicability Domain}
\author{Krishna Harish\\
\small Elkins High School, Missouri City, Texas, United States\\
\small krishnaharish2009@gmail.com}
\date{}

\begin{document}
\maketitle
\thispagestyle{empty}

This document contains: per-target Spearman correlations between each roughness descriptor and per-compound random-forest error, with the applicability-domain-controlled partial correlation and the activity-cliff fraction, for all 30 benchmark targets (Table~\ref{tab:percorr}); the full robustness sweep across neighborhood size $k$ and molecular distance metric (Table~\ref{tab:robust}); and extended graph-network architecture and training details with per-target errors (Table~\ref{tab:gnn}). All values are regenerated from the analysis pipeline by \texttt{src/make\_si\_tables.py}.

\section*{S1. Per-target correlations}

For each target, the Spearman rank correlation $\rho$ between each descriptor and the per-compound random-forest error is reported, computed on that target's test set. ``Cliff \%'' is the fraction of test compounds that participate in at least one benchmark-labeled activity cliff. ``Nbr.\ disp.\ $\mid$ AD'' is the partial Spearman correlation of the activity-free neighborhood dispersion with error after controlling (by rank residualization) for nearest-neighbor similarity and local density. Dirichlet and Lipschitz use the query's own activity; neighborhood dispersion and SALI density do not.

{\footnotesize\setlength{\tabcolsep}{4.5pt}
\begin{longtable}{l r r r r r r r}
\caption{Per-target Spearman correlations with per-compound random-forest error, and activity-cliff fractions, for all 30 MoleculeACE targets ($k = 10$, ECFP4/Tanimoto).}
\label{tab:percorr}\\
\toprule
Target & $n$ & Cliff \% & Dirichlet & Lipschitz & Nbr.\ disp. & SALI mean & Nbr.\ disp.\ $\mid$ AD \\
\midrule
\endfirsthead
\multicolumn{8}{l}{\small Table~\ref{tab:percorr} (continued)}\\
\toprule
Target & $n$ & Cliff \% & Dirichlet & Lipschitz & Nbr.\ disp. & SALI mean & Nbr.\ disp.\ $\mid$ AD \\
\midrule
\endhead
\midrule
\multicolumn{8}{r}{\small continued on next page}\\
\endfoot
\bottomrule
\endlastfoot
\siCorrRows
\end{longtable}
}

\section*{S2. Robustness across neighborhood size and distance metric}

The central correlations were recomputed for five neighborhood sizes $k$ under three distance metrics: Tanimoto similarity over ECFP4 (radius 2) and ECFP6 (radius 3) fingerprints, and Euclidean distance in a standardized 12-descriptor physicochemical space. \emph{local\_var} is a uniform-weight landscape roughness (mean squared activity gap to the $k$ neighbors, which uses the query's activity); \emph{nbr\_disp} is the activity-free neighborhood dispersion; \emph{nbr\_disp $\mid$ AD} is its partial correlation after controlling for nearest-neighbor distance and local density. Each cell is the median across the 30 targets. All fifteen $k$--metric combinations give a positive partial correlation.

\begin{table}[htbp]
\centering
\caption{Median Spearman correlation (across 30 targets) with per-compound random-forest error, across neighborhood size and distance metric.}
\label{tab:robust}
\begin{tabular}{l r r r r}
\toprule
Metric & $k$ & local\_var & nbr\_disp & nbr\_disp $\mid$ AD \\
\midrule
ECFP4 / Tanimoto & 5  & 0.641 & 0.273 & 0.226 \\
ECFP4 / Tanimoto & 10 & 0.638 & 0.275 & 0.230 \\
ECFP4 / Tanimoto & 15 & 0.628 & 0.265 & 0.207 \\
ECFP4 / Tanimoto & 20 & 0.615 & 0.240 & 0.188 \\
ECFP4 / Tanimoto & 30 & 0.587 & 0.218 & 0.186 \\
\midrule
ECFP6 / Tanimoto & 5  & 0.633 & 0.264 & 0.211 \\
ECFP6 / Tanimoto & 10 & 0.638 & 0.287 & 0.242 \\
ECFP6 / Tanimoto & 15 & 0.623 & 0.258 & 0.210 \\
ECFP6 / Tanimoto & 20 & 0.604 & 0.250 & 0.197 \\
ECFP6 / Tanimoto & 30 & 0.583 & 0.219 & 0.165 \\
\midrule
RDKit desc / Euclidean & 5  & 0.464 & 0.149 & 0.135 \\
RDKit desc / Euclidean & 10 & 0.482 & 0.163 & 0.143 \\
RDKit desc / Euclidean & 15 & 0.457 & 0.134 & 0.111 \\
RDKit desc / Euclidean & 20 & 0.460 & 0.138 & 0.115 \\
RDKit desc / Euclidean & 30 & 0.461 & 0.131 & 0.104 \\
\bottomrule
\end{tabular}
\end{table}

\section*{S3. Graph-network details and per-target errors}

\paragraph{Architecture.} The graph model is a Graph Isomorphism Network (GIN-0). Each atom carries a 16-dimensional feature vector: a one-hot atomic-number indicator over \{C, N, O, F, P, S, Cl, Br, I, other\}, degree (scaled by 4), formal charge, an aromaticity flag, total hydrogen count (scaled by 4), and a ring-membership flag. Three message-passing layers of hidden width 64 update node states as $\mathbf{H}^{(l+1)} = \mathrm{MLP}^{(l)}(\mathbf{H}^{(l)} + \mathbf{A}\,\mathbf{H}^{(l)})$, where $\mathbf{A}$ is the bond adjacency matrix and each $\mathrm{MLP}^{(l)}$ is a two-layer perceptron with ReLU activation. The graph embedding is the sum over layers of the masked mean of node embeddings, read out by a two-layer perceptron head.

\paragraph{Training.} Models are trained with Adam (learning rate $10^{-3}$) to minimize mean squared error on training-standardized targets. The \emph{fixed} regimen runs a flat schedule of 60 epochs. The \emph{tuned} regimen holds out 10\% of the training set as a validation split, trains for up to 150 epochs with learning-rate reduction on plateau and best-checkpoint early stopping (patience 20), and uses size-bucketed mini-batches of 128 to minimize padding. Per-compound test error is $|\hat{y}_i - y_i|$ in log-potency units. The tuned regimen lowers mean MAE on only 10 of 30 targets and does not change the roughness ordering reported in the main text.

\begin{longtable}{l r r r r r}
\caption{Per-target mean absolute test error (log units) for the random forest and the two graph-network regimens, with the per-target ROGI global-roughness index.}
\label{tab:gnn}\\
\toprule
Target & $n$ & RF MAE & GNN (fixed-60) & GNN (tuned) & ROGI \\
\midrule
\endfirsthead
\multicolumn{6}{l}{\small Table~\ref{tab:gnn} (continued)}\\
\toprule
Target & $n$ & RF MAE & GNN (fixed-60) & GNN (tuned) & ROGI \\
\midrule
\endhead
\midrule
\multicolumn{6}{r}{\small continued on next page}\\
\endfoot
\bottomrule
\endlastfoot
\siGnnRows
\end{longtable}

\section*{S4. Conformal prediction-interval coverage}

90\% inductive-conformal prediction intervals were built from the held-out per-compound residuals (main-text Methods). For each target the test compounds were split at random into equal calibration and evaluation halves. Coverage was made conditional on a per-compound variable in two standard ways: a locally-weighted (normalized) nonconformity score $|\hat{y}_i - y_i|/\sigma(x_i)$, and a Mondrian scheme that stratifies calibration into quintiles of the variable and calibrates within each. The conditioning variable was the random-forest tree variance, the activity-free roughness (pairwise-SALI density), the applicability domain ($1 - s_{\mathrm{NN}}$), or the combined variance-plus-roughness score. Every variant keeps the marginal coverage guarantee; they differ in where the coverage lands (Table~\ref{tab:conformal}, median across the 30 targets over 50 random splits). Conditioning on tree variance or the applicability domain leaves activity-cliff coverage at or below the unconditional level; only roughness restores it toward nominal, and the combined score recovers most of it at no width cost. The same pattern holds across the full roughness range, not only on the labeled-cliff subset: unconditional coverage falls from 0.95 in the smoothest roughness quintile to 0.87 in the roughest, and variance- and applicability-domain-conditioning fall similarly, whereas roughness (or the combined score) holds it within two points of nominal across all quintiles (main-text Figure~6A; \texttt{results/conformal\_curve.csv}).

\begin{table}[htbp]
\centering
\caption{Median 90\% conformal prediction-interval coverage and relative width across the 30 targets, by construction and conditioning variable.}
\label{tab:conformal}
\begin{tabular}{l c c c c}
\hline
Construction (conditioning variable) & Marginal & Cliff & Non-cliff & Median width (rel.) \\
\hline
Standard (unconditional) & 0.904 & 0.872 & 0.920 & 1.00 \\
Mondrian (tree variance) & 0.910 & 0.860 & 0.946 & 0.97 \\
Mondrian (applicability domain) & 0.910 & 0.861 & 0.945 & 1.01 \\
Mondrian (roughness, SALI) & 0.911 & 0.895 & 0.928 & 1.06 \\
Mondrian (variance $+$ roughness) & 0.913 & 0.876 & 0.941 & 0.99 \\
Locally-weighted (roughness, SALI) & 0.904 & 0.912 & 0.898 & 1.16 \\
Locally-weighted (applicability domain) & 0.904 & 0.820 & 0.954 & 1.14 \\
\hline
\end{tabular}
\end{table}

\end{document}
